\begin{abstract}
Power electronic converter systems (PECs) are crucial in modern power systems, particularly in the integration of renewable energy sources, electrified transportation, microgrids, and portable electronics. While PECs are known for their high energy efficiency and controllability, they encounter persistent challenges, such as nonlinear dynamics, the need for real-time control, multi-objective design trade-offs, and uncertainties in operating conditions. These challenges often lead to complex, high-dimensional, and non-convex optimization problems, which are difficult to solve using traditional model-based or rule-based techniques. Deep reinforcement learning (DRL), a model-free framework combining deep learning and reinforcement learning, emerges as a promising solution for sequential decision-making in complex environments. This paper provides a comprehensive review of DRL applications in PECs, focusing on three key areas: topology and parameter design, intelligent control, and predictive maintenance. DRL algorithms are categorized into value-based, policy-based, and actor-critic families. Their performance is evaluated based on learning efficiency, control performance, and the practical feasibility of deployment. Through critical analysis of existing approaches, this review highlights the advantages, limitations, and practical challenges of DRL techniques in PECs. It also identifies key research gaps and emerging trends, offering valuable insights for the future development of intelligent, adaptive, and autonomous converter systems. This work aims to lay the foundation for further exploration of DRL in PECs, focusing on improving system reliability and reducing operational costs in real-world applications.
\end{abstract}
